% Article fragment.  Requires amsmath, amssymb and theorem environments.
\section{Postponing the lower forest eliminates a search variable}
\label{sec:deferred-primal-forest}

The previous complete dual-forest search held the zero/one forest fixed.
The previous weighted-tree proposal alternated additive optimizations of
both forests.  For the specific objective of \emph{complete raw one/two
cancellation}, however, a different argument eliminates the lower-forest
search entirely.  The producer should first find column leaves in the full
source incidence matrix, and only then choose a zero/one forest that avoids
its pivots.  The actual geometric certificate still contracts this chosen
forest \emph{before} performing its one/two cancellations.  Only discovery
is postponed; no commutation of geometric moves is assumed.

\subsection{The graph constraint that the matrix-only model forgets}

Let $G=(V,E)$ be a finite multigraph, with loops and parallel edges allowed.
Write $c(G)$ for its number of connected components.  A spanning forest
always spans each original component.  Let $R$ be a finite row set, and let
$J$ be a disjoint set of auxiliary columns.  Consider a nonnegative integer
matrix
\[
 N=(n_{rx})_{r\in R,\ x\in E\sqcup J}.
\]
Assume that, for each row $r$, its restriction $(n_{re})_{e\in E}$ is the
literal edge-occurrence vector of a closed walk in $G$.  A union of closed
walks also suffices.  An empty restriction is permitted when a row has only
auxiliary columns.  Equivalently for the argument below, it suffices that
the uncapped occurrence vector reduce modulo two to a graph cycle:
\[
 \sum_{e\ni v} n_{re}=0\pmod 2 \qquad(v\in V),
\]
where a loop contributes twice.  One must formulate this condition before
saturation: replacing a coefficient three by two preserves singleton tests
but not parity.

A \emph{column-leaf move} chooses a live column $x$ for which
\[
 \sum_{r\in R_{\mathrm{live}}}n_{rx}=1.
\]
It deletes its unique incident row and deletes column $x$.  Other entries
are unchanged.  All earlier pivot columns have zero occurrence in every
later live row.  Thus source entries remain immutable, and the existing
sparse queue and independent timestamp checker apply unchanged.

\begin{lemma}[A closed walk cannot traverse a bridge exactly once]
\label{lem:bridge-crossing}
If $e$ is a bridge of a multigraph and $\gamma$ is a closed walk, the number
of traversals of $e$ by $\gamma$ is even.  In particular, it is never one.
\end{lemma}
\begin{proof}
Deleting $e$ partitions its original component into two vertex sets.  Every
traversal of $e$ changes the side occupied by the walk.  The initial and
final vertices agree, so the number of side changes is even.  A loop is
never a bridge.  The same proof is the mod-two cut-cycle orthogonality
identity and therefore applies to a union of closed walks.
\end{proof}

\begin{theorem}[Pivot deletion preserves the source components]
\label{thm:deferred-connectivity}
Run any finite column-leaf sequence on the full matrix $N$, and let $P$
be its pivot columns that belong to $E$.  Then $G-P$ has the same connected
components as $G$.  Consequently there is a spanning forest $T$ of $G$
with $T\cap P=\varnothing$.
\end{theorem}
\begin{proof}
Order the moves chronologically.  Auxiliary pivots do not remove graph
edges.  Suppose the next graph pivot is $e$, in row $r$, and let $P_0$ be
the graph pivots of earlier moves.  When an earlier column was a leaf, all
rows that would survive that move had zero entry in that column.  In
particular $n_{rp}=0$ for every $p\in P_0$.  The closed walk for $r$
therefore lies in $G-P_0$.  The current pivot condition gives $n_{re}=1$.
By Lemma~\ref{lem:bridge-crossing}, $e$ is not a bridge of $G-P_0$.
Deleting it preserves every component.  Induction proves the assertion.
A spanning forest of $G-P$ now also spans $G$.
\end{proof}

In graphic-matroid language, $P$ is independent in the cographic matroid:
its complement contains a spanning forest.  This observation is a
consequence of closed attaching curves, not a property of arbitrary
nonnegative matrices equipped with an unrelated graph.

\begin{theorem}[Exact removal of the primal-forest choice]
\label{thm:all-primal-forests}
Under the hypotheses above, the following conditions are equivalent:
\begin{enumerate}
\item Full-matrix column-leaf peeling deletes every row of $N$.
\item There exists a spanning forest $T$ of $G$ such that column-leaf
peeling deletes every row after columns $T$ are discarded.
\item There exists a spanning forest $T$ of $G$ such that some sequence of
raw positive-singleton cancellations deletes every row after columns $T$
are discarded, with the usual nonnegative fill update.
\end{enumerate}
Whenever these conditions hold, a full-matrix peeling sequence followed by
a spanning-forest computation in $G-P$ constructs $T$ and its complete
certificate.  There is no enumeration of primal forests.
\end{theorem}
\begin{proof}
For $(1)\Rightarrow(2)$, choose the forest supplied by
Theorem~\ref{thm:deferred-connectivity}.  No recorded pivot is in $T$.
Deleting additional nonpivot columns does not change any pivot-column sum,
so the identical row/column sequence is valid after discarding $T$.
The implication $(2)\Rightarrow(3)$ is immediate: column leaves are
zero-fill positive-singleton moves.

For $(3)\Rightarrow(2)$ use Theorem~\ref{thm:raw-completeness}: a nonnegative matrix admits a complete
positive-singleton sequence exactly when its column-leaf process deletes
all rows.  Finally, any column-leaf sequence using only columns outside
$T$ is also valid in the full matrix.  Its pivot-column sums are identical,
and the extra columns do not invalidate leaves.  Hence $(2)\Rightarrow(1)$.
\end{proof}

\begin{corollary}[A forest-independent refusal certificate]
\label{cor:all-forests-obstruction}
If full-matrix peeling leaves a nonempty stopping set $S\subseteq R$, then
no primal forest permits complete raw positive-singleton elimination.
The assertion concerns the specified geometric intersection model; it is
not a certificate that the represented manifold is not a handlebody.
\end{corollary}
\begin{proof}
Every full-matrix column sum on $S$ is zero or at least two.  Restricting
the column set by deleting any forest preserves that property.  Thus $S$
remains a stopping set in every forest-complement matrix.  The full-row
cancellation theorem excludes a complete sequence in each such matrix.
\end{proof}

The stopping certificate needs only the nonempty list of source row
identifiers.  An independent checker reconstructs the uncancelled source
occurrences, saturates each column total on that list at two, and rejects
if any total is one.  No forest search is needed to verify the universal
quantifier in the corollary.

\subsection{What is optimized when peeling stops}

Let $S$ be the terminal full-matrix column-leaf residual.  For a forest
$T$, let $S_T$ be the terminal residual after discarding $T$.  The greatest
stopping-set property gives $S\subseteq S_T$ for every $T$.  Moreover, the
forest constructed from the complete \emph{partial} full-matrix peeling
sequence attains $S_T=S$: the recorded prefix is valid, and the remaining
rows $S$ still form a stopping set.  Thus delayed forest discovery also
finds the smallest possible terminal column-leaf residual over all primal
forests.

This last statement must not be strengthened to arbitrary
positive-singleton partial schedules.  The matrix
\[
 \begin{pmatrix}1&1\\1&1\end{pmatrix}
\]
has no column leaf, but a positive singleton pivot deletes one row and
leaves the single entry two.  Both rows can be realized combinatorially as
the same closed walk $ab$ on a graph with one vertex and two loops.  The
all-forest theorem optimizes \emph{complete success}, and the partial
statement optimizes \emph{column-leaf residuals}; it does not solve a
maximum-length positive-fill cancellation problem.

\subsection{A small strict example}

Let $G$ have two vertices and three parallel edges $e_0,e_1,e_2$, all
oriented from the first vertex to the second.  Take closed walks
\[
 \gamma_0=e_0e_2^{-1}e_1e_2^{-1},\qquad
 \gamma_1=e_1e_2^{-1}.
\]
Their unsigned occurrence matrix is
\[
 N=\begin{pmatrix}1&1&2\\0&1&1\end{pmatrix}.
\]
The canonical forest $\{e_0\}$ leaves
\[
 \begin{pmatrix}1&2\\1&1\end{pmatrix},
\]
which is a stopping set.  Full-matrix peeling instead uses pivots
$(\gamma_0,e_0)$ and $(\gamma_1,e_2)$.  Their deletion leaves the spanning
forest $\{e_1\}$; contracting this forest first gives the triangular
matrix
\[
 \begin{pmatrix}1&2\\0&1\end{pmatrix},
\]
with precisely the same valid cancellation sequence.  This example is a
closed-walk incidence construction, not a claim that arbitrary graph
walks constitute authenticated manifold charts.  Actual source examples
are recorded below.

\subsection{Source rows, annular rows, and the order of authentication}

For a fixed source-authenticated two/three reduction, retain its surviving
two-handle rows and all original zero/one handles.  Each retained attaching
cycle is literally a closed walk through these handles.  The theorem
therefore applies after any chosen legal dual forest.  It eliminates the
lower forest as a search variable for each dual choice; it does not decide
which upper forest should be chosen.

For the finite bundle-boundary kernel, use the \emph{whole} retained row
set, including annular gluing rows, before choosing the source forest.
The existing midpoint tracer \texttt{source\_core\_circle} checks that
each annular core path closes in the original zero/one graph.  The
dedicated base columns are auxiliary columns $J$ and are not members of
the source forest.  A source-column pivot still has a closed core walk
through it exactly once, so Theorem~\ref{thm:deferred-connectivity}
continues to apply.  Omitting columns permanently incapable of providing a
singleton must retain the already-proved lifting hypothesis of the finite
kernel; the delayed-forest argument does not replace that hypothesis.

This extension cannot simply ignore annular rows while selecting a
forest.  A source-core leaf can cease to be a leaf after annular rows are
added.  The full finite boundary matrix must therefore be assembled first,
and its source pivot columns determine the deferred forest.

The core verifier already validates an arbitrary supplied spanning forest
against the authenticated source graph.  It accepts the new producer's
existing \texttt{normal-core-handlebody-v2} records unchanged.  For the
whole-boundary kernel we implement the new schema
\texttt{bundle-boundary-kernel-v3}: its independent checker validates the
supplied forest, reconstructs the post-forest matrix from source geometry,
and checks the same cancellation and census claims.  Older schemas retain
their canonical-forest meaning.  In every case, the geometric replay
contracts the forest before its one/two cancellations.

These certificates concern the supported fixed core or finite boundary
matrix.  Conservative returns elsewhere in the wrapper, the chosen dual
forest, omitted-column lifting conditions, and construction of a bounded
marked hierarchy successor remain separate obligations.  No assertion
about arbitrary group-word simplifications or commuting geometric moves
is used.

\subsection{Finite discovery cost}

Let $L$ be the number of literal source and annular edge occurrences,
and let $I$ be the number of nonzero saturated row/column incidences.
Constructing the matrix costs $O(L)$ indexed operations, and the column
queue uses $O(I+|R|+|E|+|J|)$ indexed work.  A depth-first or breadth-first
traversal of $G-P$ finds a spanning forest in $O(|V|+|E|)$ work.  Thus the
abstract discovery algorithm has a linear bound in these finite explicit
parameters and uses linear space.

The shipped implementation reuses the existing deterministic disjoint-set
helper, with path halving and minimum-label linking, to retain the original
component identifiers.  It does not use union by rank or size; we therefore
do not attach the usual inverse-Ackermann bound to that helper.  Sorting,
component-census loops, Python integer/dictionary operations, source
preparation, and certificate replay also have their own costs.  The
linear indexed theorem is not a claim that the entire native producer is
linear in its encoded input.  Nothing in deferred discovery requires
expanding a family represented only by an omitted genus counter.

Native core and whole-cut measurements are reported separately in
Section~\ref{sec:experiments}.
